\documentclass{tietreport}

% Import images
\usepackage{graphicx}
\usepackage{array}
\usepackage{float}

% Bibliography configuration
\usepackage[
  date=year,
  style=alphabetic,
  backend=biber,
  sorting=nty,
  maxnames=5,
  minnames=3,
  maxalphanames=4,
  minalphanames=3,
  backref=true,
  doi=false,
  isbn=false,
  url=false,
  eprint=false
]{biblatex}
\addbibresource{references.bib}

%% ==================================================
%% CUSTOMIZE THESE FIELDS FOR YOUR PROJECT
%% ==================================================

% Course/Lab header
\TitlePageHeader{%
  UCS503P - Software Engineering Project%
}

% Main project title
\ProjectTitle{%
  FinShield AI%
}

% Subtitle or project description
\ProjectSubTitle{%
  Intelligent Financial Crime Detection \& Investigation Platform%
}

% Student information (up to 4 students)
\author{%
  \texttt{1024170265} Harshit Goyal \\
  \texttt{1024170057} Arnav Gupta \\
  \texttt{1024170274} Dikshit Gupta \\
  \texttt{1024170269} Jasmandeep Sharma%
}

% Group number, degree, year, program
\TitlePageSubText{%
  Group: \texttt{1} , BE Third Year, CSE%
}

% Faculty advisor/guide name
\AdvisorName{%
   Nisha Thakur%
}

% Evaluation period and department info
\TitlePageFooterText{{%
    \large Final Stage Evaluation \\[0.2cm]
    \bfseries Computer Science and Engineering
    Department%
  } \\[0.15cm]
  Thapar Institute of Engineering and Technology,
  Patiala%
}

%% ==================================================
\let\cleardoublepage\clearpage

\begin{document}
% Generate title page
\maketitle
\tableofcontents

% ===== CHAPTER 1: INTRODUCTION =====

\chapter{Introduction}
\section{Background and Context}

Financial fraud and money laundering involve large numbers of transactions
that may contain complex and changing patterns. Traditional approaches that
depend only on fixed rules are often insufficient to identify all suspicious
activities. FinShield AI is an intelligent financial crime detection and
investigation platform that combines machine learning, rule-based analysis,
user behavior analysis, risk scoring, and investigation support.

The system receives transaction data, validates and preprocesses it,
analyzes each transaction using multiple detection approaches, and assigns
a risk level. Suspicious transactions are converted into alerts and
investigation cases for analysts. Explainable AI techniques and an AI-based
investigation assistant are integrated to support the investigation process.

\subsection{Problem Statement}

Financial institutions must quickly detect suspicious transactions without overwhelming investigators with false positives, a common pitfall of single-method systems. FinShield AI tackles this by unifying machine learning, AML rules, behavioral tracking, and hybrid risk scoring into one platform. Beyond detection, the system delivers clear risk explanations and investigative tools to help analysts act decisively.

\section{Project Objectives}

\begin{enumerate}
  \item \textbf{Objective 1:} Develop a secure platform for managing users,
    roles, and financial transaction data.
  \item \textbf{Objective 2:} Detect suspicious financial transactions using
    machine learning, AML rules, and user behavior analysis.
  \item \textbf{Objective 3:} Combine different detection results to generate
    a meaningful risk score and identify high-risk transactions.
  \item \textbf{Objective 4:} Support fraud investigation through alerts,
    explainable predictions, case management, and investigation assistance.
  \item \textbf{Objective 5:} Provide dashboards and reports for analysts,
    compliance officers, and system administrators.
\end{enumerate}

% ===== CHAPTER 2: METHODOLOGY & DESIGN =====

\chapter{Methodology and Design}

\section{System Architecture}

FinShield AI follows a modular architecture in which the frontend interacts
with the backend through REST APIs. The backend manages authentication,
transactions, fraud analysis, alerts, investigations, and communication with
the database. The fraud analysis layer contains machine learning, AML rule
processing, and user behavior analysis, whose outputs are combined by the
hybrid risk scoring component.

\subsection{High-Level Design}

The high-level design separates the system into user interaction,
application services, fraud intelligence, and data storage components.
Customers and users submit transactions through the interface. The backend
validates the transaction and passes the required information to the
fraud-analysis components. The resulting risk information is stored and
made available to analysts and compliance officers for further action.


% --------------------------------------------------
% Diagrams
% --------------------------------------------------

\begin{figure}[htbp]
  \centering
  \includegraphics[width=0.88\textwidth]{gantt_chart.jpeg}
  \caption{Three-Month Project Gantt Chart}
  \label{fig:gantt}
\end{figure}

\begin{figure}[htbp]
  \centering
  \includegraphics[width=0.88\textwidth]{ucd.jpeg}
  \caption{Use Case Diagram of FinShield AI}
  \label{fig:usecase}
\end{figure}

\begin{figure}[htbp]
  \centering
  \includegraphics[width=0.88\textwidth]{activity_diagram.jpeg}
  \caption{Activity Diagram of FinShield AI}
  \label{fig:activity}
\end{figure}

\begin{figure}[htbp]
  \centering
  \includegraphics[width=0.88\textwidth]{use_case_template_1.png}
  \caption{Use Case Template: User Authentication \& Profile Management}
  \label{fig:usecase_1}
\end{figure}

\begin{figure}[htbp]
  \centering
  \includegraphics[width=0.88\textwidth]{use_case_template_2.png}
  \caption{Use Case Template: Transaction Management}
  \label{fig:usecase_2}
\end{figure}

\begin{figure}[htbp]
  \centering
  \includegraphics[width=0.88\textwidth]{use_case_template_3.png}
  \caption{Use Case Template: ML Fraud Detection \& Risk Scoring}
  \label{fig:usecase_3}
\end{figure}

\begin{figure}[htbp]
  \centering
  \includegraphics[width=0.88\textwidth]{use_case_template_4.png}
  \caption{Use Case Template: Fraud Alert \& Notification}
  \label{fig:usecase_4}
\end{figure}

\begin{figure}[htbp]
  \centering
  \noindent\makebox[\textwidth][c]{%
    \includegraphics[width=0.88\textwidth]{use_case_template_5.png}%
  }
  \caption{Use Case Template: Fraud Case Investigation}
  \label{fig:usecase_5}
\end{figure}

\begin{figure}[htbp]
  \centering
  \includegraphics[width=0.88\textwidth]{use_case_template_6.png}
  \caption{Use Case Template: Explainable AI Analysis}
  \label{fig:usecase_6}
\end{figure}

\begin{figure}[htbp]
  \centering
  \includegraphics[width=0.88\textwidth]{use_case_template_7.png}
  \caption{Use Case Template: System \& Rule Management}
  \label{fig:usecase_7}
\end{figure}

\begin{figure}[htbp]
  \centering
  \includegraphics[width=0.88\textwidth]{use_case_template_8.png}
  \caption{Use Case Template: Risk Reports \& Compliance Analytics}
  \label{fig:usecase_8}
\end{figure}

\begin{figure}[htbp]
  \centering
  \includegraphics[width=0.88\textwidth]{collaboration_diagram.jpeg}
  \caption{Collaboration Diagram of FinShield AI}
  \label{fig:collaboration}
\end{figure}

\begin{figure}[htbp]
  \centering
  \includegraphics[width=0.88\textwidth]{sequence_diagram.jpeg}
  \caption{Sequence Diagram of FinShield AI}
  \label{fig:sequence}
\end{figure}

\begin{figure}[htbp]
  \centering
  \includegraphics[width=0.88\textwidth]{Level0_DFD.drawio.png}
  \caption{Level 0 Data Flow Diagram}
  \label{fig:dfd_0}
\end{figure}

\begin{figure}[htbp]
  \centering
  \includegraphics[width=0.88\textwidth]{Level1_DFD.drawio.png}
  \caption{Level 1 Data Flow Diagram}
  \label{fig:dfd_1}
\end{figure}

\begin{figure}[htbp]
  \centering
  \includegraphics[width=0.88\textwidth]{Level2_DFD.drawio.png}
  \caption{Level 2 Data Flow Diagram}
  \label{fig:dfd_2}
\end{figure}

\begin{figure}[htbp]
  \centering
  \includegraphics[width=0.88\textwidth]{Class_Diagram.drawio.png}
  \caption{Class Diagram}
  \label{fig:class}
\end{figure}

\subsection{Technical Stack}

\begin{itemize}
  \item \textbf{Framework/Language:} Python (FastAPI) for machine learning and
    backend services, paired with a modern web-based frontend.
  \item \textbf{Database:} PostgreSQL for structured storage of users,
    transactions, alerts, investigations, and audit records.
  \item \textbf{Tools:} Git and GitHub for version control and collaborative
    development; Jupyter Notebook for exploratory data analysis and model experiments.
  \item \textbf{Libraries:} Pandas and NumPy for data processing, Scikit-learn
    and XGBoost for model training, and SHAP for model explainability.
\end{itemize}

\section{Implementation Details}

The FinShield AI system is developed in modular layers so that each
component operates independently while adhering to a unified architecture
and common transaction data schema. Over the three-month development cycle,
the platform progressed from core transaction ingestion and baseline detection
to full integration of hybrid risk scoring, explainable AI, alert queues,
investigation case management, and monitoring dashboards.

\subsection{Implemented System Modules}

\subsubsection*{Authentication and RBAC}

The Authentication and Role-Based Access Control (RBAC) module provides
secure access to the platform. Users are authenticated before accessing
protected functionality, and permissions are enforced according to their
assigned roles: Customer/User, Analyst/Investigator, Compliance Officer,
and System Administrator. This separation ensures that each actor can
access only the functions relevant to their responsibilities.

\begin{figure}[htbp]
  \centering
  \includegraphics[width=0.88\textwidth]{authentication.jpeg}
  \caption{Authentication and RBAC}
  \label{fig:authentication_rbac}
\end{figure}


\subsubsection*{Transaction Management}

The Transaction Management module handles the submission, validation, and
storage of transaction records. When a transaction is received by the
system, it is validated against the common schema and routed through the
fraud-detection pipeline, establishing a seamless connection between user
transaction input and real-time risk screening.

\begin{figure}[htbp]
  \centering
  \includegraphics[width=0.88\textwidth]{transaction.jpeg}
  \caption{Transaction Management}
  \label{fig:transaction}
\end{figure}


\subsubsection*{ML Fraud Detection}

The ML Fraud Detection module uses engineered transaction and balance
attributes to identify potentially fraudulent activity. Development
encompassed dataset analysis, common schema definition, preprocessing,
and comparative model training. Model performance is evaluated using
imbalance-aware fraud-detection metrics—precision, recall, F1-score,
ROC-AUC, and PR-AUC—rather than relying solely on accuracy.

\begin{figure}[htbp]
  \centering
  \includegraphics[width=0.88\textwidth]{ml.jpeg}
  \caption{ML Fraud Detection}
  \label{fig:ml}
\end{figure}


\subsubsection*{AML Detection}

The AML Detection module provides rule-based and behavioral analysis for
identifying transactions indicative of money laundering. Built on the
common transaction schema, its rule engine evaluates transaction history
and feeds its risk score directly into the hybrid risk-scoring workflow
alongside the machine learning output.

\begin{figure}[htbp]
  \centering
  \includegraphics[width=0.88\textwidth]{aml.jpeg}
  \caption{AML Detection}
  \label{fig:aml}
\end{figure}


\subsection{Data and Schema Architecture}

The system enforces a common transaction-data schema across all modules.
Datasets are standardized so that the ML, AML, behavioral, backend, and
frontend components operate on a consistent representation of transaction
records. A data dictionary documents the meaning and data type of each
field, and automated preprocessing pipelines transform raw inputs before
inference.

Treating the schema as a shared contract across modules ensures clean
interoperability between the detection engines, database layer, and user
interface.

\subsection{System Integration Status}

At the final evaluation stage, all core functional layers have been
integrated into an end-to-end working application. Transaction ingestion,
ML inference, AML rule evaluation, hybrid risk scoring, SHAP explainability,
alert generation, investigation case tracking, and role-based monitoring
dashboards operate together as a cohesive platform.


% ===== CONTINUATION: IMPLEMENTATION, TESTING AND CONCLUSION =====

\subsection{Backend and API Integration}

The backend acts as the central orchestration layer between the frontend,
database, and fraud-analysis modules. Built using FastAPI, it exposes REST
API endpoints for authentication, transaction management, alerts, analytics,
and investigation cases. The backend validates incoming transaction data,
persists accepted records in PostgreSQL, and coordinates the screening
services.

The transaction-screening workflow connects the transaction-management
module with the machine-learning and AML services. The resulting risk
scores, triggered rules, and explanations are stored alongside the
transaction so that authorised staff can review them immediately.
API-based communication keeps the frontend decoupled from the internal
implementation of the detection engines.

\subsection{AML Rules, Behaviour Analysis and Risk Scoring}

The AML module evaluates transaction activity using rule-based checks and
historical account activity. Implemented rules detect unusually large
transaction amounts, high transaction frequency within a short time window,
large transfers to new receivers, and sudden spikes relative to a sender's
historical baseline. Network analysis further inspects relationships
between senders and receivers.

The AML engine outputs an AML score, a risk level, triggered rules, and
supporting observations. The hybrid risk-scoring module combines the
machine-learning (ML) fraud score and the AML score using a conditional
aggregation approach. The system first checks whether each score has
reached the threshold of 0.40.

If the ML score is below 0.40 while the AML score is at least 0.40, the
AML score is adopted as the final hybrid score. Similarly, if the AML score
is below 0.40 while the ML score is at least 0.40, the ML score is adopted
as the final hybrid score. If both scores are at least 0.40, the system
calculates a weighted score, assigning 0.60 weight to the ML score and
0.40 weight to the AML score. If both scores are below 0.40, the higher
of the two scores is selected so that the stronger signal is retained.

The scoring approach is represented as follows:

\[
H =
\begin{cases}
A, & M < 0.40 \text{ and } A \geq 0.40,\\
M, & A < 0.40 \text{ and } M \geq 0.40,\\
0.60M + 0.40A, & M \geq 0.40 \text{ and } A \geq 0.40,\\
\max(M,A), & M < 0.40 \text{ and } A < 0.40.
\end{cases}
\]

Here, $H$ represents the final hybrid risk score, $M$ represents the
machine-learning fraud score, and $A$ represents the AML score. A score
exactly equal to 0.40 is treated as having reached the threshold.

The final hybrid score is classified into three risk levels: HIGH for
scores greater than or equal to 0.70, MEDIUM for scores greater than or
equal to 0.40 and less than 0.70, and LOW for scores below 0.40. This
conditional formulation prevents a low score from one engine from diluting
a strong fraud or laundering signal detected by the other.

\subsection{Alerts and Investigation Support}

Transactions flagged as medium or high risk are automatically routed to
the alert queue for review by authorised analysts. Each alert displays
the risk level, hybrid score, triggered rules, and review status. The
case-management module enables investigators to open formal cases,
record investigation notes, update case statuses, and log final outcomes,
completing the workflow from detection to resolution.

\subsection{Explainability and Monitoring}

FinShield AI incorporates explainability features to make model outputs
transparent and actionable. SHAP-based feature contributions highlight
which transaction attributes most strongly influenced a machine-learning
prediction. In addition, the platform provides analytical dashboards for
monitoring screened transactions, active alerts, and ongoing investigation
cases, empowering analysts to make informed decisions rather than relying
on opaque risk scores.

\subsection{Component Diagram}

The component diagram illustrates the main software components of FinShield AI
and the interfaces between the frontend, backend services, fraud analysis
modules, and database.

\begin{figure}[htbp]
  \centering
  \IfFileExists{component_diagram.jpeg}{
    \includegraphics[width=0.95\textwidth]{component_diagram.jpeg}
  }{
    \fbox{\parbox[c][5cm][c]{0.88\textwidth}{\centering
      \textbf{Component Diagram Placeholder}\\[0.3cm]
      Upload \texttt{component\_diagram.jpeg} to the Overleaf project.}}
  }
  \caption{Component Diagram of FinShield AI}
  \label{fig:component}
\end{figure}

\subsection{Deployment Diagram}

The deployment diagram illustrates the runtime architecture of the
frontend, backend API, machine-learning and AML services, and PostgreSQL
database, including the communication protocols between these nodes.

\begin{figure}[htbp]
  \centering
  \IfFileExists{deployment_diagram.jpeg}{
    \includegraphics[width=0.95\textwidth]{deployment_diagram.jpeg}
  }{
    \fbox{\parbox[c][5cm][c]{0.88\textwidth}{\centering
      \textbf{Deployment Diagram Placeholder}\\[0.3cm]
      Upload \texttt{deployment\_diagram.jpeg} to the Overleaf project.}}
  }
  \caption{Deployment Diagram of FinShield AI}
  \label{fig:deployment}
\end{figure}

% ===== CHAPTER 3: RESULTS AND EVALUATION =====

\chapter{Results and Evaluation}

\section{Testing Methodology}

Testing was carried out across multiple levels to verify the correctness of
individual modules and the end-to-end interaction between the frontend,
backend, database, and fraud-analysis services.

\begin{itemize}
  \item \textbf{Unit Testing:} Individual AML rules and machine-learning
  inference functions were tested with normal and suspicious transaction
  inputs.
  \item \textbf{Integration Testing:} API and backend integration tests
  verified that transaction submission, risk scoring, database persistence,
  and retrieval operate reliably together.
  \item \textbf{Functional Testing:} End-to-end test scenarios verified
  transaction screening, risk categorisation, alert generation, and routine
  transaction clearance.
  \item \textbf{User Interface Testing:} Frontend forms, dashboards, and
  alert views were verified to confirm smooth transaction submission and
  result inspection.
  \item \textbf{Hybrid Scoring Tests:} Conditional aggregation logic was
  verified across boundary conditions where one score is below 0.40, both
  scores are at least 0.40, and both scores are below 0.40.
\end{itemize}

Table~\ref{tab:hybridtests} presents test cases verifying the conditional
hybrid scoring logic across all decision branches.

\begin{table}[htbp]
  \centering
  \small
  \begin{tabular}{cccc}
    \hline
    \textbf{ML Score} & \textbf{AML Score} &
    \textbf{Expected Hybrid Score} & \textbf{Scoring Condition} \\
    \hline
    0.30 & 0.80 & 0.80 & AML score selected \\
    0.80 & 0.30 & 0.80 & ML score selected \\
    0.60 & 0.80 & 0.68 & Weighted calculation \\
    0.90 & 0.70 & 0.82 & Weighted calculation \\
    0.20 & 0.30 & 0.30 & Higher score selected \\
    \hline
  \end{tabular}
  \caption{Verification Test Cases for Conditional Hybrid Risk Scoring}
  \label{tab:hybridtests}
\end{table}

\section{Selected Functional Test Cases}

Four representative functional test cases were executed on the completed
system to validate high-, medium-, and low-risk transaction screening
workflows.

\subsection{Test Case 1.1: Large Transfer Emptying the Sender Account}

A transfer of 200,000 is submitted from sender account \texttt{s2}, whose
balance decreases from 200,000 to zero. The receiver account is
\texttt{r2}, with its balance increasing from 5,000 to 205,000. The
receiver has no earlier transaction history in this test setup.

\textbf{Observed outcome:} The transaction is classified as high risk
due to the large transfer amount, account depletion, and transfer to a
new receiver. An alert is generated in the queue for analyst review.

\begin{figure}[htbp]
  \centering
  \IfFileExists{test_case_1_1.png}{
    \includegraphics[width=0.92\textwidth]{test_case_1_1.png}
  }{
    \fbox{\parbox[c][4.2cm][c]{0.88\textwidth}{\centering
      \textbf{Test Case 1.1 Screenshot Placeholder}\\
      Upload \texttt{test\_case\_1\_1.png} to Overleaf.}}
  }
  \caption{Execution Screenshot: Test Case 1.1}
  \label{fig:test11}
\end{figure}

\subsection{Test Case 1.2: Large Transfer to a New Receiver After Small Transfers}

Two earlier transfers of 1,000 and 2,000 are first submitted by sender
\texttt{S05} to receivers \texttt{R05\_OLD1} and \texttt{R05\_OLD2}.
A new transfer of 60,000 is then submitted by the same sender to
\texttt{R05\_NEW}. The sender balance changes from 97,000 to 37,000.

\textbf{Observed outcome:} The transaction is classified as medium risk
because it represents a large transfer to a previously unseen receiver
and deviates sharply from the sender's earlier transaction amounts. An
alert is generated accordingly.

\begin{figure}[htbp]
  \centering
  \IfFileExists{test_case_1_2_a.png}{
    \includegraphics[width=0.92\textwidth]{test_case_1_2_a.png}
  }{
    \fbox{\parbox[c][4.2cm][c]{0.88\textwidth}{\centering
      \textbf{Test Case 1.2 (Part 1) Screenshot Placeholder}\\
      Upload \texttt{test\_case\_1\_2\_a.png} to Overleaf.}}
  }
  \caption{Execution Screenshot: Test Case 1.2 (Transaction Submission \& History)}
  \label{fig:test12a}
\end{figure}

\begin{figure}[htbp]
  \centering
  \IfFileExists{test_case_1_2_b.png}{
    \includegraphics[width=0.92\textwidth]{test_case_1_2_b.png}
  }{
    \fbox{\parbox[c][4.2cm][c]{0.88\textwidth}{\centering
      \textbf{Test Case 1.2 (Part 2) Screenshot Placeholder}\\
      Upload \texttt{test\_case\_1\_2\_b.png} to Overleaf.}}
  }
  \caption{Execution Screenshot: Test Case 1.2 (Risk Scoring \& Alert Result)}
  \label{fig:test12b}
\end{figure}

\subsection{Test Case 1.3: Burst of Transfers to Multiple Receivers}

Five transfers of 1,000 each are submitted by sender \texttt{S14} to five
different receivers, \texttt{R14\_A} through \texttt{R14\_E}. A sixth
transfer of 150,000 is then submitted to the new receiver
\texttt{R14\_NEW}. The sender balance changes from 295,000 to 145,000
for the final transaction.

\textbf{Observed outcome:} The final transaction is classified as high
risk due to the large transfer amount combined with the sender's rapid
burst of transfers across multiple receivers. An alert is created for
investigation.

\begin{figure}[htbp]
  \centering
  \IfFileExists{test_case_1_3.png}{
    \includegraphics[width=0.92\textwidth]{test_case_1_3.png}
  }{
    \fbox{\parbox[c][4.2cm][c]{0.88\textwidth}{\centering
      \textbf{Test Case 1.3 Screenshot Placeholder}\\
      Upload \texttt{test\_case\_1\_3.png} to Overleaf.}}
  }
  \caption{Execution Screenshot: Test Case 1.3}
  \label{fig:test13}
\end{figure}

\subsection{Test Case 1.4: Small Routine Payment}

A routine payment of 500 is submitted from sender \texttt{s1} to receiver
\texttt{r1}. The sender balance changes from 10,000 to 9,500 and the
receiver balance changes from 5,000 to 5,500. No earlier transactions
are supplied for either account in this test setup.

\textbf{Observed outcome:} The transaction receives a low-risk
classification. No AML rules are triggered and no alert is generated.

\begin{figure}[htbp]
  \centering
  \IfFileExists{test_case_1_4.png}{
    \includegraphics[width=0.92\textwidth]{test_case_1_4.png}
  }{
    \fbox{\parbox[c][4.2cm][c]{0.88\textwidth}{\centering
      \textbf{Test Case 1.4 Screenshot Placeholder}\\
      Upload \texttt{test\_case\_1\_4.png} to Overleaf.}}
  }
  \caption{Execution Screenshot: Test Case 1.4}
  \label{fig:test14}
\end{figure}

\section{Performance Metrics}

The machine-learning models were evaluated using precision, recall,
F1-score, ROC-AUC, and PR-AUC. These metrics are critical for fraud
detection because fraudulent transactions represent a small fraction of
overall volume, where accuracy alone would mask poor minority-class
detection.

The holdout evaluation results on the PaySim dataset are summarised in
Table~\ref{tab:modelmetrics}. These results reflect benchmark dataset
performance rather than live banking traffic.

\begin{table}[htbp]
  \centering
  \small
  \begin{tabular}{lccccc}
    \hline
    \textbf{Model} & \textbf{Precision} & \textbf{Recall} &
    \textbf{F1-score} & \textbf{ROC-AUC} & \textbf{PR-AUC} \\
    \hline
    Logistic Regression & 0.0337 & 0.9811 & 0.0652 & 0.9962 & 0.6188 \\
    Random Forest       & 1.0000 & 0.9976 & 0.9988 & 0.9991 & 0.9981 \\
    XGBoost             & 0.8755 & 0.9970 & 0.9323 & 0.9994 & 0.9941 \\
    \hline
  \end{tabular}
  \caption{Holdout Evaluation Results on the PaySim Dataset}
  \label{tab:modelmetrics}
\end{table}

The results demonstrate strong holdout performance for the tree-based
ensemble models. Feature attribution analysis indicates that the models
rely heavily on sender balance-related fields, which are strong indicators
in the PaySim dataset. Additional time-based validation and cross-dataset
testing remain important before deploying models on live banking streams.

\section{Analysis}

The implemented system successfully integrates transaction screening with
machine-learning inference, AML rule evaluation, hybrid risk scoring, and
case management. Across the four functional test scenarios, the platform
accurately differentiated high-risk account-draining and burst transfers,
medium-risk behavioral anomalies, and routine low-risk payments.

In model evaluation, Random Forest and XGBoost achieved high precision
and recall on the holdout split, whereas Logistic Regression exhibited
high recall but low precision due to excessive false positives. Combining
tree-based ML predictions with deterministic AML rules provides balanced
coverage across both known statistical fraud patterns and explicit
compliance thresholds.

Certain operational limitations remain. The current evaluation relies on
benchmark datasets and simulated transaction flows rather than live core
banking feeds. Rule thresholds and hybrid weights are configurable system
parameters that should be calibrated with institutional domain experts,
and generated risk scores serve as decision-support signals for human
investigators rather than automated legal determinations.

% ===== CHAPTER 4: CONCLUSION =====

\chapter{Conclusion}

\section{Summary of Work}

FinShield AI is an end-to-end intelligent platform for financial
transaction screening and crime investigation support. The system
integrates a web frontend, FastAPI backend services, PostgreSQL storage,
machine-learning fraud detection, and rule-based AML analysis. It
delivers a complete operational workflow spanning conditional hybrid risk
scoring, SHAP explainability, alert triaging, investigation case
management, and analytical monitoring.

\section{Key Findings}

\begin{itemize}
  \item Combining machine-learning predictions with AML rules and
  transaction-history analysis provides complementary coverage of both
  statistical anomalies and explicit compliance violations.
  \item Precision, recall, F1-score, ROC-AUC, and PR-AUC provide a far
  more reliable evaluation of imbalanced fraud models than accuracy alone.
  \item Historical sender and receiver context is essential for detecting
  behavioral deviations, new-payee risks, and multi-recipient transfer bursts.
  \item Conditional hybrid scoring effectively preserves a strong risk
  signal when one detection engine scores below the 0.40 threshold, while
  blending both signals when both detect elevated risk.
  \item Explainable feature attributions and structured case workflows
  significantly assist human investigators without replacing expert
  compliance judgment.
\end{itemize}

\section{Future Work and Recommendations}

The following enhancements are recommended for future iterations of the platform:

\begin{enumerate}
  \item Incorporate graph neural networks (GNNs) to uncover deeper
  multi-hop money laundering rings across account networks.
  \item Integrate real-time streaming ingestion (such as Apache Kafka)
  for high-throughput live banking transaction processing.
  \item Expand model validation using out-of-time splits, drift
  monitoring, and multi-institution datasets.
  \item Enhance automated regulatory reporting (such as SAR generation)
  and expand natural-language investigation summaries.
  \item Conduct formal penetration testing, encryption hardening, and
  load testing prior to production deployment.
\end{enumerate}

\section{Final Remarks}

FinShield AI demonstrates how modern software engineering principles,
machine learning, and rule-based compliance engines can be unified into an
effective financial crime detection and investigation platform. By pairing
hybrid risk scoring with transparent explanations and structured case
management, the system provides a practical foundation for accelerating
fraud and AML investigations.

% ===== BIBLIOGRAPHY =====
\printbibliography

\end{document}